\documentclass[a4paper,11pt]{article}
\usepackage[margin=27mm]{geometry}
\usepackage{fontspec,amsmath,amssymb,amsthm}
\setmainfont{STIXGeneral}
\setlength{\emergencystretch}{2em}
\usepackage{unicode-math}\setmathfont{STIX Two Math}
\usepackage[hidelinks]{hyperref}
\newcommand{\Sym}{\operatorname{Sym}}
\newcommand{\gen}[1]{\langle #1\rangle}
\newtheorem{theorem}{Theorem}\newtheorem{lemma}[theorem]{Lemma}\newtheorem{corollary}[theorem]{Corollary}
\title{Gluing permutation markings through a common point}
\author{Human Compression project}\date{Private author edition; independent review pending}
\begin{document}\maketitle
\begin{abstract}
We give a sufficient condition for identifying one point in two permutation
C-markings while preserving every subgroup intersection. The condition asks
that each generated subgroup contain every permutation supported on the orbit
of the distinguished point. A pure leaf transposition guarantees it.
The proof uses the resulting direct factors and avoids a rank induction.
It gives the first gluing theorem of Leemans--Mulpas as a corollary
under the literal graph hypotheses verified below. The rest of that paper,
including its second construction and counterexample families, is not yet
recovered by this draft.
\end{abstract}

\section{Permutations supported on the distinguished orbit}
All sets in this note are finite. If $R\subseteq\Omega$, then $\Sym(R)$
means the subgroup of $\Sym(\Omega)$ fixing $\Omega\setminus R$ pointwise.
If $H\leq\Sym(\Omega)$ preserves $R$, write $H^{\Omega\setminus R}$ for
its induced permutation group on the complement, extended by the identity on
$R$. This is an actual subgroup of $\Sym(\Omega)$, but it need not be a
subgroup of $H$ without a further hypothesis.

\begin{lemma}\label{lem:factor}
Let $a\in\Omega$, let $H\leq\Sym(\Omega)$, and let $R=a^H$.
If $\Sym(R)\leq H$, then
\[
 H=\Sym(R)\times H^{\Omega\setminus R}.
\]
\end{lemma}
\begin{proof}
Every element of $H$ preserves its orbit $R$. Given $h\in H$, extend its
restriction to $R$ by the identity outside $R$. This extension belongs to
$\Sym(R)\leq H$. Cancelling it from $h$ leaves exactly the extension of
the complementary restriction. Consequently the entire complementary image
is a subgroup of $H$. The two factors have disjoint supports, commute, and
intersect trivially, proving the formula.
\end{proof}

A family $(s_i)_{i\in I}$ of distinct nonidentity involutions is a
\emph{C-marking} if, on writing $H_J=\gen{s_j:j\in J}$, it satisfies
\[
 H_J\cap H_K=H_{J\cap K}\qquad(J,K\subseteq I).
\]
The group is the group generated by the family. No string condition is
included in this definition. The intersection property alone does not force
intersections of point orbits to agree: the two generators $(12)$ and
$(12)(34)$ give a C-marking of a group of order four, but both singleton
generator subgroups have orbit $\{1,2\}$ at $1$, whereas their intersection
is trivial. Condition~\eqref{eq:support} below rules out precisely the
missing supported transposition needed in our argument.

\begin{lemma}\label{lem:orbits}
Suppose a C-marking on $\Omega$ has the property
\begin{equation}\label{eq:support}
 \Sym(a^{H_J})\leq H_J\qquad(J\subseteq I).
\end{equation}
Then its distinguished orbits satisfy
\[
 a^{H_J}\cap a^{H_K}=a^{H_{J\cap K}}.
\]
\end{lemma}
\begin{proof}
One inclusion follows from subgroup containment. If $b$ belongs to the
left side and $b\ne a$, the transposition $(a\ b)$ belongs to both
$H_J$ and $H_K$ by \eqref{eq:support}. The intersection property puts
it in $H_{J\cap K}$, so $b$ belongs to the required orbit. The case
$b=a$ follows from the identity element.
\end{proof}

\section{The complete intersection theorem}
Let $\Omega_1\cap\Omega_2=\{a\}$ and put
$\Omega=\Omega_1\cup\Omega_2$. Permutations on either set are extended
by the identity on the other set minus $a$. The index sets of the two
markings are understood to be disjoint.

\begin{theorem}\label{thm:gluing}
For $\nu=1,2$, let $(s_{\nu i})_{i\in I_\nu}$ be a C-marking on
$\Omega_\nu$ satisfying \eqref{eq:support} at $a$.
The union of the two markings is a C-marking on $\Omega$.
More precisely, for $J_\nu\subseteq I_\nu$, put
\[
 H_{\nu,J_\nu}=\gen{s_{\nu i}:i\in J_\nu},\qquad
 R_{\nu,J_\nu}=a^{H_{\nu,J_\nu}},\qquad
 K_{\nu,J_\nu}=H_{\nu,J_\nu}^{\Omega_\nu\setminus R_{\nu,J_\nu}}.
\]
Then the subgroup on the selected marks is
\begin{equation}\label{eq:joined}
 G_{J_1,J_2}
 =\Sym(R_{1,J_1}\cup R_{2,J_2})
       \times K_{1,J_1}\times K_{2,J_2}.
\end{equation}
In particular the resulting marking again satisfies \eqref{eq:support}
at $a$, and any finite number of such markings can be glued at $a$.
\end{theorem}
\begin{proof}
Lemma~\ref{lem:factor} gives the decomposition of each selected input
subgroup. Two symmetric groups on sets meeting in $a$ generate the symmetric
group on their union: each contains all transpositions $(a\ x)$ in its own
set, and those transpositions generate the symmetric group of the union.
This includes a singleton set, whose symmetric group is trivial.
All complementary factors have disjoint supports from this union and from
each other. This proves \eqref{eq:joined}. It also identifies the orbit of
$a$ in the joined group and proves its asserted property.

We prove the intersection property for two arbitrary choices
$J=(J_1,J_2)$ and $L=(L_1,L_2)$. Abbreviate the symmetric-factor supports
in \eqref{eq:joined} by $R_J$ and $R_L$. By Lemma~\ref{lem:orbits} in
each input,
\[
 R_J\cap R_L
   =R_{1,J_1\cap L_1}\cup R_{2,J_2\cap L_2}
   =R_{J\cap L}.
\]
Let $g\in G_J\cap G_L$. It preserves both $R_J$ and $R_L$, hence their
intersection $R_{J\cap L}$. Formula~\eqref{eq:joined} for the common labels
contains all permutations supported on this intersection. Cancel from $g$
its restriction there; call the remaining permutation $g'$. Thus
$g'\in G_J\cap G_L$ and $g'$ fixes $R_{J\cap L}$ pointwise, including $a$.

An element of $G_J$ can send a point from
$\Omega_1\setminus\{a\}$ to $\Omega_2\setminus\{a\}$ only inside
$R_J$: its other two factors each have support on one side. The same
holds for $G_L$. Hence any such crossing by $g'$ would start in
$R_J\cap R_L$, which $g'$ fixes. The inverse gives the reverse direction.
Therefore $g'$ preserves each $\Omega_\nu$ and fixes $a$.
Its restriction to $\Omega_\nu$ lies in both $H_{\nu,J_\nu}$ and
$H_{\nu,L_\nu}$: in formula~\eqref{eq:joined} its restriction on the
symmetric-factor support lies in $\Sym(R_{\nu,J_\nu})$, and the
remaining restriction belongs to $K_{\nu,J_\nu}$, with the identical
argument for $L$. Input intersection now puts these restrictions in
$H_{\nu,J_\nu\cap L_\nu}$. Their extensions fix $a$ and have disjoint
remaining supports, so their product $g'$ belongs to $G_{J\cap L}$.
Restoring the cancelled permutation proves $g\in G_{J\cap L}$.
The reverse inclusion is immediate.

All marks remain involutions. Distinct input marks cannot become equal
across the two sides: a common permutation would be supported on
$\Omega_1\cap\Omega_2=\{a\}$ and hence would be the identity.
Thus the enlarged family is a C-marking. Induction now proves the final
assertion, since the property at $a$ has already been preserved.
\end{proof}

\section{Trees of attached configurations}
The hypothesis at other possible attachment points survives the construction.
This is useful when a marking is assembled from several pieces, rather than
from a single pair.

\begin{lemma}\label{lem:ports}
In Theorem~\ref{thm:gluing}, suppose $b\in\Omega_1$ and
$\Sym(b^{H_{1,J}})\leq H_{1,J}$ for every $J\subseteq I_1$.
The glued marking has the same property at $b$.
\end{lemma}
\begin{proof}
Fix arbitrary selected labels on both sides. If $b$ belongs to the selected
orbit of $a$ in the first input, its orbit in the joined group is $R_J$,
and formula~\eqref{eq:joined} supplies its supported symmetric group.
Otherwise the orbit of $b$ stays inside the complementary first-side factor:
its orbit is exactly its former orbit, because that factor has the full induced
action of $H_{1,J_1}$ there. Its supported symmetric group was already in
$H_{1,J_1}$ and remains in the joined group. These two cases exhaust $b$.
\end{proof}

\begin{corollary}\label{cor:tree}
Let $T$ be a nonempty finite tree. At each vertex $v$ place a C-marking of rank $r_v$
on a finite set $\Omega_v$ of size $n_v$. Initially all these sets and all
marking index sets are disjoint. For each incident tree edge choose a different
point of $\Omega_v$, and require the supported-orbit condition
\eqref{eq:support} at every chosen point, for every subset of that marking.
For each edge of $T$, identify its two chosen points. The union of all markings
on the resulting set is a C-marking of degree $N$ and rank $r$, with
\[
 N=\sum_v n_v-|E(T)|,\qquad r=\sum_v r_v,\qquad
 N-r=1+\sum_v(n_v-r_v-1).
\]
It retains the supported-orbit condition at every point that had this
condition in one of the original pieces.
\end{corollary}
\begin{proof}
A one-vertex tree requires no gluing. Otherwise remove a leaf. The smaller
tree has the asserted C-marking by induction. Its attachment point retains
the required condition by Lemma~\ref{lem:ports}. The removed piece meets the
assembled set in exactly that one point: the original sets were disjoint,
ports inside a piece were distinct, and the incidence pattern has no cycle.
Theorem~\ref{thm:gluing} and Lemma~\ref{lem:ports} complete the induction.
Each edge identifies exactly one pair of previously distinct points, while
no labels are identified. This proves the first two formulas. The third uses
$|E(T)|=|V(T)|-1$.
\end{proof}

For a path of pieces, assume each input marking is string. Require each
left attachment point, when present,
to be moved only by the first mark of its piece, and each right attachment
point, when present, only by its last mark. Require these endpoint marks
to be pure transpositions, and the two points in any internal piece to be
distinct. A rank-one piece may use its single transposition at both ends.
Order the pieces along the path and their marks compatibly. The union is a string C-marking: within each piece the
string relations hold, and the only cross-piece pairs with intersecting
supports are the two endpoint marks at a joint, which are consecutive.
The pure leaf lemma below verifies all hypotheses of Corollary~\ref{cor:tree}.
A rank-one piece on two points contributes zero to the sum $N-r-1$.
Thus arbitrarily long chains of such pieces preserve degree minus rank.
This is a construction theorem; it does not assert that every marking of a
fixed degree minus rank has such a decomposition or a bounded remaining core.

\section{Pure leaf transpositions}
\begin{lemma}\label{lem:leaf}
Suppose a marking on $\Omega$ has a distinguished mark
$s_0=(a\ b)$, and every other mark fixes $a$.
Then every subgroup on a subset of its marks satisfies
\eqref{eq:support}, whether or not the marking has the intersection property.
\end{lemma}
\begin{proof}
If the subset omits $s_0$, its orbit of $a$ is the singleton $\{a\}$,
and there is nothing to prove. Otherwise let $M$ be the subgroup generated
by the other selected marks and put $O=b^M$. Since $M$ fixes $a$, it
preserves $O\subseteq\Omega\setminus\{a\}$. Conjugating $(a\ b)$ by
$M$ gives every transposition $(a\ x)$ with $x\in O$; these generate
$\Sym(\{a\}\cup O)$. The selected subgroup preserves this set and is
transitive on it, so it is precisely the orbit of $a$. The required
supported symmetric group has just been exhibited inside it.
\end{proof}

\begin{corollary}\label{cor:string}
Take two string C-markings, each with first mark a pure transposition
$(a\ b_\nu)$ and all later marks fixing $a$. Identify their distinguished
points $a$ and no other points. Reverse the order of the first marking
and concatenate it with the second. The resulting ordered marking is a
string C-marking.
\end{corollary}
\begin{proof}
Lemmas~\ref{lem:leaf} and Theorem~\ref{thm:gluing} prove every subset
intersection. Within each old marking the string commutations persist.
Across the two sides, only their distinguished first marks can move the
common point. Those two marks are consecutive in the new order.
Every other cross pair has disjoint supports and therefore commutes.
\end{proof}

If an input graph with this pure leaf is connected, its group is the full
symmetric group: removing the leaf leaves the orbit $O$ of the preceding
proof equal to all remaining points. For a disconnected input the component
containing the leaf is still a supported symmetric direct factor, and the
induced action on the other components is the complementary factor of
Lemma~\ref{lem:factor}. No primitivity theorem is required for either claim.

\begin{lemma}\label{lem:bridge}
Let $\Omega=A\sqcup B$, let $L\leq\Sym(A)$ and $M\leq\Sym(B)$,
and let $a\in A$, $b\in B$. In the group
$H=\gen{L,M,(a\ b)}$, the orbit $R$ containing $a,b$ is
$a^L\cup b^M$, and $\Sym(R)\leq H$. Consequently
$H=\Sym(R)\times H^{\Omega\setminus R}$.
\end{lemma}
\begin{proof}
Both $L$ and $M$ preserve the union $R=a^L\cup b^M$, as does $(a\ b)$.
Conjugating $(a\ b)$ by the commuting groups $L$ and $M$ gives every
transposition $(x\ y)$ with $x\in a^L$ and $y\in b^M$.
These are the edges of a connected complete bipartite graph, and generate
$\Sym(R)$. For completeness, its edges generate every transposition on $R$:
along a path, conjugating an already obtained transposition by the next
edge transposition extends its endpoint by one step. Hence all transpositions
from one fixed vertex to the others lie in the generated group, and they
generate the symmetric group. This also proves that $R$ is one orbit.
Lemma~\ref{lem:factor} gives the final decomposition.
\end{proof}

\section{Exact relation to the existing source and limitations}
Leemans--Mulpas, \emph{Two gluing methods for string C-group representations
of the symmetric groups}, arXiv:2504.17535v1, states its first gluing operation
in Theorem 1.1 (physical pages 1--2). The displayed root has only the first
edge and the remaining graph has labels at least one. Under that literal
hypothesis the first mark is a pure transposition and all other marks fix the
root; Corollary~\ref{cor:string} proves the conclusion. Its edge-addition
Lemma 3.2 is the case in which one input has one mark. The supported-factor
argument also replaces the primitive-group step in its Lemma 3.1 when the
single bridge separates disjoint lower and higher label supports. All lower labels act on one side and all higher labels on the other;
Lemma~\ref{lem:bridge} therefore applies verbatim, including disconnected
components. The source uses perfect-split primitivity and its primitive-group
transposition lemma for this step; neither is needed in this replacement.

These are the source's conclusions. The stronger arbitrary C-marking
criterion in Theorem~\ref{thm:gluing} is a proposed common formulation;
no historical priority or complete-source replacement is claimed here.
In particular a labelled edge belonging to an involution with additional
transpositions is not a pure transposition. The source itself gives a failure
of the weakened construction (physical page 6):
\[
 (12)(34),\ (23),\ (34)
\]
is a string C-marking on four points, but adjoining $(15)$ at its first
end fails the intersection property. If the new order is
\[
 t_{-1}=(15),\quad t_0=(12)(34),\quad t_1=(23),\quad t_2=(34),
\]
then $\gen{t_0,t_1,t_2}=S_{\{1,2,3,4\}}$ has order $24$,
$\gen{t_{-1},t_0,t_1}=S_5$, and
$\gen{t_0,t_1}$ has order $8$. Their intersection is therefore strictly
larger than the subgroup on the common labels. Indeed, the first subgroup
contains $(34)$ and hence $(12)=t_0(34)$, giving the three adjacent
transpositions of $S_4$. The group $\langle t_0,t_1\rangle$ is dihedral of
order eight since $t_0t_1$ is a four-cycle. It is transitive on the first four
points and fixes $5$, so conjugating $(15)$ by it yields every $(j5)$,
$1\leq j\leq4$; these generate $S_5$. This recorded source
counterexample is not a new disproof.

Neither the one-point intersection nor the tree assumption can be discarded.
For example, the two tree markings $(13),(23)$ and $(14),(24)$ on sets
meeting in $\{1,2\}$ have disjoint labels, but their groups intersect in
$\langle(12)\rangle$. Likewise three rank-one markings on two points,
attached cyclically, yield $(12),(23),(31)$; the first two already generate
the third, so the intersection property fails. These examples explain the
set-intersection restrictions in the general construction.

The full later gluing families, their auxiliary identifications and the
source's uniform counterexample to a proposed second operation remain to be
recovered. Likewise this theorem alone gives no classification of all markings
with fixed degree-minus-rank: a separate theorem must find suitable splits
or bound the exceptional part. No numerical comparison or new source credit
is attached to this draft.

The accompanying Lean file verifies six statements about actual permutation
subgroups: membership and monotonicity of point orbits, the intersection of
two orbits under the supported-transposition hypotheses, its C-family
consequence, and the conjugation and orbit-inclusion steps for a pure leaf.
It allows arbitrary underlying types. The restriction decompositions,
complete gluing proof and tree construction remain written proofs. The
GAP program checks twelve paired and two iterated finite gluings, as well as
four failed weakenings; these
finite cases are diagnostics, not premises of the quantified results.
\end{document}
